% !TeX root = ../偏微分方程笔记.tex
% 由主文档通过 \input 载入；请编译主文档。
\setcounter{chapter}{3} % 保留原书章号。
\chapter{二阶线性方程的分类与标准形}
\label{chap:classification}

最高阶导数决定方程对局部高频变化的主要响应，因而称为\emph{主部（principal part）}。二维二阶方程的主部与一个二次型对应，分类首先是线性代数问题。

\section{主部、主符号与二维分类}
考虑实系数方程
\begin{equation}
 A u_{xx}+2B u_{xy}+C u_{yy}+D u_x+E u_y+F u=G,
 \label{eq:second-order-general}
\end{equation}
其主部矩阵为 $M=\bigl(\begin{smallmatrix}A&B\\B&C\end{smallmatrix}\bigr)$。在 $M\ne0$ 的点，若 $B^2-AC<0$，两特征值同号，称\emph{椭圆型（elliptic）}；若 $B^2-AC>0$，两者异号，称\emph{双曲型（hyperbolic）}；若 $B^2-AC=0$，矩阵秩为一，通常称二维主部为\emph{抛物型（parabolic）}。$M=0$ 时二阶主部消失，应单独视为退化点。

$u_{xx}+u_{yy}=0$、$u_{tt}-c^2u_{xx}=0$ 和 $u_t-\kappa u_{xx}=0$ 分别给出三类典型模型。但秩一的主部本身不保证方程具有热传导那样的时间演化结构，低阶项仍需检查。

判别式中的 $B$ 是混合导数系数的一半。例如 $u_{xx}+u_{xy}+u_{yy}=0$ 对应 $B=1/2$，而不是 $B=1$。从二次型看，椭圆型在任意非零方向上都有同号曲率权重，双曲型则存在正负两类方向；热方程把时间视为独立变量时，主部只有空间二阶导数，演化方向由一阶时间项补充。因此“方程二阶”与“需要两个初始条件”并不等价，后者要看时间导数的阶数与问题类型。

\section{坐标变换与标准形}
设 $\xi=\xi(x,y)$、$\eta=\eta(x,y)$ 为 $C^2$ 可逆局部坐标，写 $u(x,y)=v(\xi,\eta)$。以 $u_{xx}$ 为例，
\[
 u_{xx}=v_{\xi\xi}\xi_x^2+2v_{\xi\eta}\xi_x\eta_x+v_{\eta\eta}\eta_x^2
         +v_\xi\xi_{xx}+v_\eta\eta_{xx}.
\]
若 $J=\partial(\xi,\eta)/\partial(x,y)$，新主部矩阵为 $JMJ^{\mathsf T}$。由于 $\det(JMJ^{\mathsf T})=(\det J)^2\det M$，分类在非退化坐标变换下不变。

常系数时可先用正交变换对角化 $M$，再伸缩坐标，将主部化为 $v_{\xi\xi}+v_{\eta\eta}$、$v_{\xi\xi}-v_{\eta\eta}$ 或 $v_{\xi\xi}$，必要时整体乘以非零常数。变系数情形逐点对角化不等于已经找到全局坐标；坐标的可积性和雅可比非零均须另行处理。

\begin{example}[直接用特征坐标消去平方项]
把 $u_{xx}-4u_{xy}+3u_{yy}=0$ 化为标准形，并写出全部局部 $C^2$ 解。
\end{example}
\begin{solve}
主部可因式分解为 $(\partial_x-\partial_y)(\partial_x-3\partial_y)$。取 $\xi=y+x$、$\eta=y+3x$，则雅可比为 $-2\ne0$，且
\[
 \partial_x=\partial_\xi+3\partial_\eta,\qquad
 \partial_y=\partial_\xi+\partial_\eta.
\]
两个因子分别化为 $2\partial_\eta$ 与 $-2\partial_\xi$，故 $-4v_{\xi\eta}=0$。在变换后的小矩形内积分得到 $v=P(\xi)+Q(\eta)$，即 $u=P(y+x)+Q(y+3x)$。这里 $P,Q$ 是任意 $C^2$ 函数；若有边界或初始数据，还需用它们确定这两个函数，不能把含任意函数的表达式当作定解问题的最终答案。
\end{solve}

\section{特征曲线与传播结构}
设曲线局部为 $\varphi(x,y)=\text{常数}$，其法向为 $\nabla\varphi$。当
\[
 A\varphi_x^2+2B\varphi_x\varphi_y+C\varphi_y^2=0
\]
时称为特征曲线。用切向微分表示即
\[
 A(dy)^2-2B\,dx\,dy+C(dx)^2=0.
\]
双曲型有两族实特征方向，椭圆型没有非零实特征方向，秩一抛物型有一个重方向。波动方程在 $(x,t)$ 平面中给出 $dx=\pm c\,dt$。取 $\xi=x-ct$、$\eta=x+ct$，方程变成 $v_{\xi\eta}=0$，这是达朗贝尔公式的起点。

\section{高维主符号与分类的边界}
对实对称矩阵 $M(x)$，定义\emph{主符号（principal symbol）}$p(x,\zeta)=\zeta^{\mathsf T}M(x)\zeta$。若存在 $\theta>0$，对所有 $\zeta\in\mathbb R^n$ 和所考虑的 $x$ 都有
\[
 \theta|\zeta|^2\le\zeta^{\mathsf T}M(x)\zeta\le\theta^{-1}|\zeta|^2,
\]
则在此正号约定下称主部一致椭圆。负定主部可先整体换号。高维双曲或抛物演化方程还需指定时间方向，不能使用一个二维判别式概括。

特里科米方程（Tricomi equation）$y u_{xx}+u_{yy}=0$ 在 $y>0$ 为椭圆型，在 $y<0$ 为双曲型，在 $y=0$ 退化为秩一主部。此例说明方程类型可以随位置改变。

\section{例题与习题}
\begin{example}
化简 $u_{xx}+4u_{xy}+u_{yy}=0$。
\end{example}
\begin{solve}
$B=2$，故 $B^2-AC=3>0$。令 $\xi=(x+y)/\sqrt2$、$\eta=(x-y)/\sqrt2$，雅可比为 $-1$，主部变为 $3v_{\xi\xi}-v_{\eta\eta}$。再令 $r=\xi/\sqrt3$，得到 $w_{rr}-w_{\eta\eta}=0$。所有附加条件也必须按相同变换重写。
\end{solve}
\begin{enumerate}
 \item 对 $u_{xx}+2u_{xy}+u_{yy}+u_x=0$ 取 $\xi=x+y$、$\eta=x-y$，验证所得方程为 $4v_{\xi\xi}+v_\xi+v_\eta=0$。
 \item 写出 $y u_{xx}+u_{yy}=0$ 在 $y<0$ 内的特征曲线微分方程。
 \item 给出一个处处椭圆但在整个无界区域上不一致椭圆的系数矩阵。
\end{enumerate}

